% BHU PYQ -- Manually transcribed & error-corrected
% Paper: BCH-215 : Fundamentals of Business Finance
% Exam: B.Com. (Hons.) Semester III Examination, 2023-24

\documentclass[11pt,a4paper]{article}
\usepackage[a4paper,top=2.5cm,bottom=2.5cm,left=2.8cm,right=2.8cm,headheight=14pt]{geometry}
\usepackage{amsmath,amssymb}
\usepackage[T1]{fontenc}
\usepackage[utf8]{inputenc}
\usepackage{lmodern,microtype,fancyhdr,enumitem,hyperref,lastpage,parskip}

\hypersetup{
    colorlinks=true,
    urlcolor=black,
    linkcolor=black,
    pdftitle={BCH-215: Fundamentals of Business Finance -- BHU B.Com. (Hons.) Sem III 2023-24}
}

\pagestyle{fancy}
\fancyhf{}
\renewcommand{\headrulewidth}{0.4pt}
\renewcommand{\footrulewidth}{0.4pt}
\fancyhead[L]{\small   \textit{Banaras Hindu University}}
\fancyhead[R]{\small   \textit{BCH-215: Fundamentals of Business Finance}}
\fancyfoot[C]{\small Page  \thepage\ of \pageref{LastPage}}


\newlist{parts}{enumerate}{2}
\setlist[parts,1]{label=(\roman*),leftmargin=*,itemsep=5pt,topsep=3pt}
\setlist[parts,2]{label=(\alph*),leftmargin=*,itemsep=3pt,topsep=2pt}

\newcommand{\pts}[1]{\hfill   \textit{[#1]}}


\begin{document}


\begin{center}
  {\large\bfseries BANARAS HINDU UNIVERSITY}\\[2pt]
  {\normalsize B.Com. (Hons.) --- Semester III Examination, 2023-24}\\[6pt]
  \rule{0.8\linewidth}{0.5pt}\\[6pt]
  {\large\bfseries Paper No. BCH-215: Fundamentals of Business Finance}\\[4pt]
  {\normalsize Time: 3 Hours\hfill Maximum Marks: 70}
\end{center}

\rule{\linewidth}{0.4pt}

\smallskip



\noindent   \textit{Instructions: Answer all questions. All questions carry equal marks (14 marks each).}

\bigskip




\noindent   \textbf{Question 1.}
Discuss clearly the meaning and scope of finance function. ``The financial man should not only be a money man but also a business man.'' Elucidate the statement.\pts{14}

\centerline{   \textbf{OR}}

Explain the concept and rationale of wealth maximization objective. How far is wealth maximization objective free from the ideological and operational objections raised against profit maximization?\pts{14}

\medskip\rule{\linewidth}{0.2pt}




\noindent   \textbf{Question 2.}
``Financial planning is the key to the successful business operations.'' Explain and discuss the basic characteristics of financial plan of a joint stock company.\pts{14}

\centerline{   \textbf{OR}}

Explain the meaning of financial plan. What are the characteristics of a sound financial plan? Briefly discuss the factors affecting such plans.\pts{14}

\medskip\rule{\linewidth}{0.2pt}




\noindent   \textbf{Question 3.}
What objectives tests would you employ to ascertain the state of under capitalization? Analyse its causes and effects.\pts{14}

\centerline{   \textbf{OR}}


\begin{parts}
  \item ``Over-capitalisation has nothing to do with redundance of capital in an enterprise.'' Critically discuss this statement and show how to find whether or not an enterprise is over capitalized or under-capitalised. Briefly state the effects of over capitalization and under-capitalisation.
  \item The Balance Sheet of S. Co. Ltd., Agra as on 30th June, 2022 is as follows:
  
  
\begin{center}
  \small
  
\begin{tabular}{|l|r||l|r|}
  \hline
   \textbf{Liabilities} &   \textbf{Amount (Rs.)} &   \textbf{Assets} &   \textbf{Amount (Rs.)} \\
  \hline
  Share Capital: & & Sundry Assets & 5,00,000 \\
  \quad 4,000 Shares of Rs. 100 each & 4,00,000 & & \\
  Reserve and Surplus & 60,000 & & \\
  Current Liabilities & 40,000 & & \\
  \hline
   \textbf{Total} &   \textbf{5,00,000} &   \textbf{Total} &   \textbf{5,00,000} \\
  \hline
  \end{tabular}
  \end{center}
  
  If the average profit of the company is Rs. 40,000 and similar companies earn a profit of 10\% on capital employed, state whether this company suffers from over-capitalisation or not.
\end{parts}\pts{14}

\medskip\rule{\linewidth}{0.2pt}




\noindent   \textbf{Question 4.}
Mention the meaning, merits and demerits of equity and preference shares as a source of long-term finance for a company.\pts{14}

\centerline{   \textbf{OR}}

``Debentures occupy a very important place in corporate finance.'' Discuss this statement and point out the limitations of financing through debentures.\pts{14}

\medskip\rule{\linewidth}{0.2pt}




\noindent   \textbf{Question 5.}
``Unfortunately money-suppliers, accountants, traders and economists hold divergent views about working capital.'' Elucidate this statement and explain the kinds and sources of working capital.\pts{14}

\centerline{   \textbf{OR}}

Write a note on the analysis of working capital. What is the operating cycle method of working capital forecasting?\pts{14}

\vfill
\rule{\linewidth}{0.4pt}

\begin{center}\small
\href{https://bhu-exam.vercel.app/}{Click here to visit our website}\\[4pt]
\href{https://me-aryan.vercel.app/}{Click here to visit our contributor}\\[4pt]
\href{https://quashan-me.vercel.app/}{Click here to visit our contributor}
\end{center}

\end{document}
